\subsection{Concurrency Model and Support} \label{concurrency-model} % done

FreeBSD's concurrency is centered around its \textbf{1:1 threading model}. This model denotes that for every user-level
thread, it is backed by a corresponding kernel thread, which has its own kernel stack \parencite{mckusick2014_process}. The
kernel creates the illusion of concurrency by scheduling system resources among the set of ready processes
\parencite{mckusick2014_process}. Multiple threads of the same or different processes can be executed in true parallel fashion
on a multiprocessor, which has different CPUs \parencite{mckusick2014_process}. Those processors share a single instance of
the kernel code and data. This parallelism is enforced by the ULE scheduler with per-CPU run queues and topology-aware
affinity decisions, as described in \ref{load-balancing-distribution} \parencite{freebsd_smp_handbook}. An interrupt
that arrives on a CPU than the one running the kernel code is forwarded using \textit{inter-processor interrupts (IPIs)}
\parencite{freebsd_smp_handbook}. This concurrency model was not initially adopted by FreeBSD. On its 5.0 release, it
incorporates \textbf{Kernel Scheduled Entities (KSE)}, which is an intermediary model that adopted an M:N threading
model \parencite{freebsd_libthr}. The KSE multiplexed user threads over a small number of kernel-level virtual processors,
which became the default in FreeBSD 5.3 \parencite{freebsd_libthr}. It was eventually replaced by the now 1:1 threading
model in FreeBSD 7.0 since it outperformed the M:N model enforced by the \texttt{libkse} library \parencite{freebsd_libthr}.
In FreeBSD 8.0, KSE was eventually removed \parencite{freebsd_libthr}. While the M:N model reduces kernel-thread overheads,
it actually introduced more complexity in terms of I/O blocking and signal handling \parencite{freebsd_libthr}, which
might be one of the key reasons as to why it was replaced with the now 1:1 model. Synchronizing concurrent processes
is achieved by FreeBSD's adapive blocking mutex primitive \parencite{freebsd_locking}. If the owner of a contended mutex
is running on another CPU, the acquiring thread does not yield but rather spins \parencite{freebsd_locking}. Inversely,
if the owner is not running, the thread blocks until the mutex is released \parencite{freebsd_locking}. This mechanism
results in short hold times and achieves a good balance between spin and sleep locking \parencite{freebsd_locking}.

NetBSD also adopted a 1:1 threading model which was first introduced in NetBSD 5.0 \parencite{rasiukevicius2009}.
Similar to FreeBSD's concurrency model, each user-level thread, called a \texttt{pthread}, is mapped to a kernel-level
process or the \textbf{light-weight process (LWP)} \parencite{rasiukevicius2009}. Inside the kernel, traditional processes
and POSIX threads are both represented as LWPs which are being scheduled similarly \parencite{rasiukevicius2009}.
Similar to FreeBSD, this current threading model replaced an older hybrid M:N model which was first proposed by
\textcite{anderson1991}, which was implemented by \textcite{williams_scheduler} for the same reasons. At the same time,
NetBSD rewrote its entire synchronization subsystem. Previously in NetBSD 4.0, the kernel used to rely on spinlocks
and a per-CPU interrupt priority level which is set with \texttt{spl(9)} \parencite{netbsd_locking}. These mechanisms
did not support kernel preemption, hence they were not suitable for multiprocessors \parencite{netbsd_locking}. Hence,
they were replaced with thread-based adaptive mutexes, reader-writer locks, and condition variables, which were adopted
in NetBSD 5.0 \parencite{netbsd_locking}.

The NetBSD scheduler treats every LWP uniformly and supports the following POSIX scheduling policies:

\begin{itemize}
	\item \texttt{SCHED\_OTHER}: the default time-sharing policy, in which each LWP carries a dynamically recalculated
	      priority and time quantum \parencite{rasiukevicius2009}.

	\item \texttt{SCHED\_FIFO}: a real-time first-in-first-out policy whose threads hold fixed priorities and run until
	      they complete or block \parencite{rasiukevicius2009}.

	\item \texttt{SCHED\_RR}: a real-time round-robin policy in which threads at the same priority are cycled
	      within their class \parencite{rasiukevicius2009}.
\end{itemize}

\noindent
Additionally, the kernel maintains 224 priority levels \parencite{sh4nnu2020}. Of these 224 levels, 64 are available
to user-level scheduling policies \parencite{sh4nnu2020}. The Common Scheduler Framework, which was introduced in
NetBSD 5.0, lets the 4BSD and M@ scheduling algorithms previously discussed in Section \ref{resource-allocation-strats}
be selected at compile time \parencite{sh4nnu2020}.

The two systems reached the same destination by different routes. Both abandoned M:N threading once they judged its
conceptual elegance not worth its implementation complexity \parencite{freebsd_libthr}. FreeBSD reached that
verdict by retiring KSE, and NetBSD reached it by retiring scheduler activations \parencite{rasiukevicius2009}. Neither
system now multiplexes user threads onto a separate pool of kernel contexts, so neither pays the overhead that motivated
the M:N designs in the first place. What still separates them is how they implemented the 1:1 model. In FreeBSD,
each user thread is bound to a distinct kernel thread, which keeps its model close to what most kernel code already assumes
\parencite{mckusick2014_process}. NetBSD instead treats processes and threads alike as LWPs \parencite{rasiukevicius2009}. That
single abstraction is what lets the NetBSD scheduler apply the same time-sharing and real-time policies to every entity,
exactly as the policy list above reflects \parencite{rasiukevicius2009}.

\subsection{Process Synchronization} \label{process-synchronization} % done

The FreeBSD kernel offers a variety of synchronization primitives, each tuned to a specific access pattern. Table
\ref{freebsd-sync-primitives} summarizes them. The primary primitive is the \textit{adaptive blocking mutex}, which
protects short-lived shared data \parencite{freebsd_locking}. As previously mentioned in Section \ref{concurrency-model},
the mutex spins if the holder is on-CPU and blocks otherwise \parencite{freebsd_locking}. It also fully supports
priority propagation through the turnstile mechanism to prevent priority inversion \parencite{freebsd_locking}.
This primitive prevents a thread from sleeping while holding a thread. This is enforced by \texttt{INVARIANTS},
which simplifies lock-order reasoning and prevents mutexes from spanning long operations such as disk I/O operations
\parencite{freebsd_locking}. \textit{Reader-writer locks} or \texttt{rwlock()} extend this with shared and exclusive
modes \parencite{freebsd_locking}. They propagate priority for exclusive writers but not for anonymous shared readers
\parencite{freebsd_locking}. They also may not be held across an unbounded sleep \parencite{freebsd_locking}. The
\textit{read-mostly lock} or \texttt{rmlock()} targets rarely-written data \parencite{freebsd_rmlock}. It resolves
the anonymous-reader problem by propagating priority to both readers and writers through a caller-supplied
\texttt{rm\_priotracker} struct \parencite{freebsd_rmlock}. This choice places the allocation burden on the caller's
stack and also maintains a per-CPU reader list, which can become a bottleneck in it of itself when many read-mostly
locks are being used at the same time \parencite{freebsd_rmlock}.

\begin{table}[!htb]
	\centering
	\begin{tabular}{|c|p{0.72\textwidth}|}
		\hline
		\textbf{Primitive} & \textbf{Description}                                                                                                                                                        \\
		\hline\hline
		\texttt{mutex()}   & Adaptive blocking lock for short-lived shared data. Spins if the holder is on-CPU and blocks otherwise. Propagates priority via turnstiles. May not be held while sleeping. \\ \hline
		\texttt{rwlock()}  & Reader-writer lock permitting many concurrent readers or one writer. Propagates priority for writers but not anonymous readers. May not be held across unbounded sleep.     \\ \hline
		\texttt{rmlock()}  & Read-mostly lock for rarely-written data. Propagates priority to readers and writers via a caller-supplied \texttt{rm\_priotracker}. Tracks readers in a per-CPU list.      \\ \hline
		\texttt{sx()}      & Shared/exclusive lock that may be held during unbounded sleep, suiting VFS and buffer-cache operations. Does not support priority propagation.                              \\ \hline
		\texttt{sema()}    & Counting semaphore for resource pools and producer-consumer patterns. Has no owner concept. Largely deprecated in favor of mutexes and condition variables.                 \\ \hline
		\texttt{condvar()} & Condition variable paired with a blocking lock. Atomically releases the lock and sleeps, then reacquires it on wakeup to avoid lost wakeups.                                \\ \hline
	\end{tabular}
	\caption{FreeBSD Kernel Synchronization Primitives \parencites{freebsd_locking}{freebsd_rmlock}{freebsd_sema}}
	\label{freebsd-sync-primitives}
\end{table}

The remaining FreeBSD primitives trade priority propagation for sleep tolerance or generality. The \textit{counting semaphore} or
\texttt{sema()} synchronizes access to a pool of resources \parencite{freebsd_sema}. It also supports producer-consumer
patterns where one thread acquires and another releases \parencite{freebsd_sema}. \textcite{freebsd_sema} deprecates them
because of their complexity and inefficiency when compared to mutexes and condition variables. The \textit{shared/exclusive lock}
or \texttt{sx()} resembles a reader-writer lock but may be held during an unbounded sleep \parencite{freebsd_locking}.
This makes it suitable for VFS and buffer-cache operations that block waiting threads for I/O \parencite{freebsd_locking}.
However, unlike mutexes and reader-writer locks, it does not propagate priority, so a high-priority thread waiting on an
\texttt{sw} lock cannot accelerate its holder \parencite{freebsd_locking}. \textit{Condition variables} or \texttt{condvar()}
pair with any blocking lock and let a thread sleep atomically \parencite{freebsd_locking}. They release the associated
lock before yielding and reacquire it on wakeup, which prevents lost wakeups \parencite{freebsd_locking}.

NetBSD's synchronization primitive set is also similar to FreeBSD's, but it is largely documented around explicit
usage preconditions and execution contexts. The adaptive mutex or \texttt{mutex()} requires a constant acquisition
order across all code paths \parencite{netbsd_mutex}. \textcite{netbsd_mutex} prescribes an optimistic \texttt{mutex\_tryenter()}
pattern where reverse-order acquisition is unavoidable. If the attempt fails, both locks are released and reacquired in
the original order \parencite{netbsd_mutex}. This pattern is a code-level compensation for the absence of a witness-like
detector in production kernels \parencite{netbsd_mutex}. The reader-writer lock or \texttt{rwlock()} is documented with
three explicit preconditions for appropriate use according to \textcite{netbsd_rwlock}. For instance, the data must be read
far more often than written on the order of thousands of cycles \parencite{netbsd_rwlock}. Furthermore, strong synchronization
semantics are required for this approach \parencite{netbsd_rwlock}. When these do not hold, \textcite{netbsd_rwlock}
recommends restructuring the implementation that needs fewer or weaker synchronization points. For read-mostly data, NetBSD
also offers a passive serialization mechanism with \texttt{pserialize()} for lock-free reads \parencite{netbsd_pserialize}.
Readers bracket a critical path with \texttt{pserialize\_read\_enter()} and \texttt{pserailize\_read\_exit()}, while
a writer calls \texttt{pserialize\_perform()} to block until all concurrent readers have exited \parencite{netbsd_pserialize}.
This permits reads even in software-interrupt context at \texttt{IPL\_SOFTCLOCK} without acquiring any lock
\parencite{netbsd_pserialize}. The cost is that \texttt{pserialize\_perform()} is confined to the thread context
\parencite{netbsd_pserialize}. Table \ref{netbsd-context-matrix} lists which primitives are valid in each execution
context. The matrix aids kernel developers in selecting the appropriate primitive for a given context, which is a direct
expression of NetBSD's focus on correctness and portability. For userland inter-process synchronization, NetBSD supports
POSIX and System V IPC \parencite{netbsd_tuning}. This includes named semaphores via \texttt{sem\_open()}, SysV semaphore
arrays, shared memory through \texttt{shmget}/\texttt{shmat}, and message queues \parencite{netbsd_tuning}. These are all
inspectable by calling \texttt{ipcs()} \parencite{netbsd_tuning}. SysV shared memory can optionally be pinned in physical
memory via \texttt{kern.ipc.shm\_use\_phys} to keep it from being paged out, which can benefit database workloads
\parencite{netbsd_tuning}.

\begin{table}[!htb]
	\centering
	\begin{tabular}{|p{0.34\textwidth}|c|c|c|}
		\hline
		\textbf{Primitive}                            & \textbf{Thread} & \textbf{Soft Interrupt} & \textbf{Hard Interrupt} \\
		\hline\hline
		Adaptive mutex (\texttt{mutex()})             & Yes             & Context-dependent       & Context-dependent       \\ \hline
		Reader-writer lock (\texttt{rwlock()})        & Yes             & Yes                     & No                      \\ \hline
		Condition variable (\texttt{condvar()})       & Yes             & Partial                 & No                      \\ \hline
		Passive serialization (\texttt{pserialize()}) & Yes             & Yes                     & No                      \\ \hline
		Interrupt priority level (\texttt{spl()})     & Yes             & No                      & No                      \\ \hline
		Atomic operations (\texttt{atomic\_ops()})    & Yes             & Yes                     & Yes                     \\ \hline
	\end{tabular}
	\caption{Validity of NetBSD Synchronization Primitives by Execution Context \parencites{netbsd_locking}{netbsd_atomic_ops}}
	\label{netbsd-context-matrix}
\end{table}

While both FreeBSD and NetBSD adopts the same synchronization primitives, they are documented differently by their
manual pages. Kernel developers working in FreeBSD need to explicitly select primitives that provide a good balance
of priority propagation and sleep tolerance for a specific workload. FreeBSD also offered a rich variety of lock types
to choose from \parencite{freebsd_locking}. This is unlike NetBSD, that offers a much smaller set which is being governed
by explicit preconditions. This smaller set nudges NetBSD developers towards minimal synchronization rather than
more specialized synchronizationThis smaller set nudges NetBSD developers towards minimal synchronization rather than
more specialized synchronization. The documentation also provides a context-validity matrix that better represents which
contexts are ideal for specific named locks. Another point of contrast is FreeBSD's \texttt{rmlock} and NetBSD's
\texttt{pserialize}. While both make reads cheap on rarely-written data, the \texttt{rmlock} performs priority propagation
while \texttt{pserialize} drops locking on the read side. FreeBSD's approach therefore aligns with its emphasis on
correct behavior under contention, while NetBSD's implementation keeps synchronization small but forces
blocks on writes until all the readers have been freed.

\subsection{Inter-Process Communication (IPC)} \label{ipc} % done

While FreeBSD provides a variety of IPC mechanisms, two of them are most used, which are pipes and sockets. \textit{Pipes}
are unidirectional and unnamed byte-stream channels between processes that share a common parent \parencite{bellthomas2020}.
\textit{Sockets} are bidirectional and more adaptable with support for more complex packet structures and network
communication \parencite{bellthomas2020}. Notably, FreeBSD no longer stacks pipes on top of sockets but rather adopts
a separate implementation that is built directly on the VM system \parencite{bellthomas2020}. A study conducted by
\textcite{bellthomas2020} on FreeBSD's performance found that pipes outperform local socket pairs by 63\% on average
across all configurations. This gap comes from how pipes implement VM optimizations, which transfer data between processes
by page-flipping instead of copying \parencite{bellthomas2020}. This cannot be replicated by sockets since their versatile
design requires an in-kernel buffer model \parencite{bellthomas2020}. Optimal throughput for both pipes and sockets can
reach around 32--64K range with transmission buffers, which results in increased buffer-management overhead and performance
degradation when exceeding that range \parencite{bellthomas2020}.

In sending files over the network, FreeBSD provides the \texttt{sendfile()} system call, which first appeared in FreeBSD 3.0
\parencite{freebsd_sendfile}. \texttt{sendfile()} performs zero-copy file-to-socket transfer by building an \texttt{mbuf}
chain that points to resident file pages in the cache \parencite{freebsd_sendfile}. This chain is then passed to the TCP/IP stack
\parencite{freebsd_sendfile}. This eliminates two copy operations and the memory pressure of duplicating cached data
\parencite{freebsd_sendfile}. This approach significantly contributes to FreeBSD's strength in managing high-volume content.
However, this is only operable in the disk-to-network direction and cannot be served as a general zero-copy proxy
\parencite{freebsd_sendfile}. FreeBSD supports the full System V IPC suite alongside POSIX variants at the highest
bandwidth length \parencite{tyk_sysvipc}. SysV shared memory eliminates copying entirely by mapping the same physical pages
into multiple address spaces \parencite{tyk_sysvipc}.  This makes it the dominant mechanism for database
workloads \parencite{tyk_sysvipc}. However, this mechanism does not provide built-in synchronization, hence processes
must coordinate access with separate semaphores or mutexes \parencite{tyk_sysvipc}. Inside the kernel, the shared/exclusive
lock plays a role similar to the IPC by serializing access to the VFS layer and buffer cache while permitting concurrent reads
\parencite{freebsd_smp_lock_strategies}. FreeBSD often layers lock types as well, protecting the parent-process pointer with both the
\texttt{proctree\_lock} \texttt{sx} lock and a per-process mutex \parencite{freebsd_smp_lock_strategies}.

\subsection{Handling Critical Sections and Race Conditions} \label{critical-sections} % done

Critical sections in FreeBSD can be entered with \texttt{critical\_enter()} and exited with \texttt{critical\_exit()}
\parencite{mckusick2014_process}. It prevents context switches on the current CPU, so neither preemption nor migration can occur
while inside it \parencite{mckusick2014_process}. Critical sections carry less overhead than locks and allow for short operations on
per-CPU run queues and allocators \parencite{mckusick2014_process}. However, they protect only the current CPU and does not provide
for system-wide shared data \parencite{mckusick2014_process}. This is significant because FreeBSD's preemptive kernel allows a thread
to migrate to another CPU at any time an interrupt triggers preemption, which can make unprotected per-CPU access unsafe
\parencite{freebsd_smp_handbook}. A critical section defers any ongoing preemption until \texttt{critical\_exit()} is called
\parencite{freebsd_smp_handbook}. Consequently, it can blocks interrupt threads on the current processor, which can raise interrupt
latency \parencite{freebsd_smp_handbook}. For finer-grained ordering without a full lock, FreeBSD provides atomic operations with
acquire and release semantics \parencite{freebsd_atomic}. On one hand, an acquire barrier on lock acquisition prevents loads and
stores from being reordered before it \parencite{freebsd_atomic}. On the other hand, a release barrier on unlock prevents
reordering after it \parencite{freebsd_atomic}. Both of these barriers provide the memory ordering that makes writing data inside a
critical section visible to other CPUs that will later acquire the same lock \parencite{freebsd_atomic}.
FreeBSD exposes improperly handled critical sections and race conditions during development through the \texttt{INVARIANTS} kernel
option \parencite{freebsd_debugging}. This option enables thousands of runtime integrity checks converting silent data corruption
into a kernel panic with a stack trace \parencite{freebsd_debugging}. Because of its runtime cost, \texttt{INVARIANTS} is
disabled in production \parencite{freebsd_debugging}. Instead, the kernel relies on its locking hierarchy together with the \texttt{WITNESS}
verifier described in Section \ref{handling-deadlock} to expose violations before deployment \parencite{freebsd_debugging}.

NetBSD handles kernel preemption with \texttt{kpreempt\_disable()} and \texttt{kpreempt\_enable()}
\parencite{netbsd_kpreempt}. Preemption is also disabled in several cases, which includes cases whenever the interrupt priority
level is raised above \texttt{IPL\_NONE} through \texttt{spl()}, whenever a spinning mutex is held, or whenever \texttt{kernel\_lock}
is held \parencite{netbsd_kpreempt}. Each of these prevents the calling thread from migrating during a sensitive operation
\parencite{netbsd_kpreempt}. NetBSD offers the architecture-independent \texttt{atomic\_ops()} read-modify-write interface for
fine-grained ordering \parencite{netbsd_atomic_ops}. Its guarantees are deliberately weaker than on a single architecture
\parencite{netbsd_atomic_ops}. Where hardware compare-and-swap is available, stores become globally visible before the atomic
operation completes \parencite{netbsd_atomic_ops}. Code needing surrounding loads and stores ordered must add \texttt{membar\_ops()}
memory barriers since atomic operations are only ordered with one another \parencite{netbsd_atomic_ops}. This requirement adds
extra burden when compared with x86-only systems where atomics imply full barriers \parencite{netbsd_atomic_ops}.
This mechanism coexists with the legacy \texttt{spl()} interface, which is  still valid in thread context but not in interrupt
context \parencite{netbsd_locking}. This masks interrupts to keep handlers from racing code that holds an elevated IPL
\parencite{netbsd_locking}.

The two systems share a prevention-first stance, but they differ in where they place the burden and how they verify it.
Both disable preemption or migration to protect per-CPU data. Both also supply architecture-aware atomic and barrier primitives
NetBSD adds the \texttt{mutex\_tryenter()} reverse-order recovery pattern as a code-level guard
against lock-order inversion \parencite{netbsd_mutex}. A notable difference is how systems detect race conditions.
FreeBSD's \texttt{INVARIANTS} and \texttt{WITNESS} turn race conditions and ordering errors into immediate panics
accessible in development environments \parencite{freebsd_debugging}. For NetBSD, it relies more on convention and the
context matrix. It also lacks a production-grade detector that is comparable in performance with FreeBSD. The NetBSD jails work
discussed in Section \ref{resource-contention} showed that \texttt{secmodel\_eval} calls from hot-path-adjacent subsystems
could interact with lock ordering and recursion in UVM and scheduling paths to produce a full and unrecoverable freeze
\parencite{petermann_jails}.

\subsection{Deadlock Revisited} \label{deadlock-revisited} % done

Building on the prevention-by-ordering strategies introduced in Section \ref{handling-deadlock}, both systems guard against
particular structural deadlock classes that arise specifically from concurrency. FreeBSD addresses a case in which a thread holds a
spin mutex, is interrupted, and the interrupt handler tries to acquire the same lock, which would wedge both parties
\parencite{freebsd_mutex}. It prevents this by requiring spin mutexes to disable interrupts on the local CPU while held
\parencite{freebsd_mutex}. This is so that no handler can contest the lock there \parencite{freebsd_mutex}.
A second case arises when a thread holding a spin mutex is preempted and the new thread spins forever waiting for a lock the
evicted holder cannot release \parencite{freebsd_smp_handbook}. FreeBSD treats spin-lock hold as a preemption-critical region
and disables preemption through the critical-section mechanism \parencite{freebsd_smp_handbook}. This prevents the holder from
being evicted before releasing \parencite{freebsd_smp_handbook}. Potential ordering violations can be caught during the development
of the kernel by the \texttt{WITNESS} option \parencite{freebsd_kerneldebug}. It maintains a directed graph of lock-acquisition
orders and checks for cycles at every acquire \parencite{freebsd_kerneldebug}. It also prints a warning and stack trace on detection
\parencite{freebsd_kerneldebug}. It can drop into the debugger under \texttt{WITNESS\_KDB} and can skip the heavily-exercised spin
locks under \texttt{WITNESS\_SKIPSPIN} \parencite{freebsd_kerneldebug}. Its limitation, as previously mentioned in Section
\ref{handling-deadlock}, is that it is disabled in production, so violations confined to rare paths may surface only after
deployment \parencite{freebsd_kerneldebug}.

Beyond deadlock, FreeBSD must contend with two hazards related to liveness. The first is starvation. ULE bounds the maximum latency a
thread of a given priority experiences by the product of its slice size and the number of runnable threads at that priority
\parencite{mckusick2014_process}. Also, its two-queue timeshare structure works to keep higher-priority batch jobs from starving
lower-priority ones \parencite{mckusick2014_process}. Even so, as mentioned in Section \ref{resource-performance},
\textcite{bouron2018} found that ULE can starve background threads even within a single application. While that behavior can benefit
the interactive threads, it breaks fairness assumptions for workloads with identical threads \parencite{bouron2018}. The
second hazard is livelock. BSD-derived systems are susceptible to \textit{receive livelock}, in which the system spends
all its time servicing receives interrupts under extreme network load and delivers no packets to applications \parencite{mogul1996}.
This is because interrupt handling has absolute priority \parencite{mogul1996}. The solution which was developed for 4.4BSD
is to schedule network interrupt processing as deliberately as process execution, which can be done either through per-interface
rate limiting or by switching to polling when load is extreme \parencite{mogul1996}.

NetBSD prevents deadlock through the same convention-based ordering with the addition of the \texttt{LOCKDEBUG} diagnostic option
mentioned in Section \ref{handling-deadlock}. A complementary safeguard is the context-validity matrix in Table
\ref{netbsd-context-matrix} \parencite{netbsd_locking}. It specifies that adaptive mutexes are invalid in hard interrupt context,
that is \texttt{rwlock()} is valid in thread and soft interrupt context but not hard interrupt context \parencite{netbsd_locking}.
Furthermore, \texttt{spl()} is valid only in thread context \parencite{netbsd_locking}. The matrix also consider cross-context
deadlocks such as taking a sleeping lock from an interrupt handler and misuse panics in debug builds \parencite{netbsd_locking}.
For starvation, NetBSD's time-sharing queue boosts I/O-bound LWPs on wakeup and dynamically recalculates each LWP's priority and
quantum, which cycles threads of the same priority in round-robin fashion \parencite{rasiukevicius2009}. Real-time
\texttt{SCHED\_FIFO} threads can by design starve lower-priority threads indefinitely, which is the reason as to why only the
superuser may assign \texttt{SCHED\_FIFO} \parencite{rasiukevicius2009}. For livelock, the jails freeze mentioned in Section
\ref{resource-contention} is NetBSD's documented real-world case \parencite{petermann_jails}. It is the absence of a
production-ready deadlock or livelock detector implying that such conditions can presently be found
only through exhaustive testing \parencite{petermann_jails}.

FreeBSD pairs convention with the mature \texttt{WITNESS} verifier \parencite{freebsd_kerneldebug}. It then
adds structural defenses, which include interrupt-disabling spin mutexes, preemption-disabled spin regions, and an
explicit polling escape from receive livelock \parencite{freebsd_smp_handbook}. Each named concurrency hazard thus has a
concrete answer in FreeBSD. While NetBSD relies on convention, the context matrix, and compile-time \texttt{LOCKDEBUG}
\parencite{netbsd_deadlock_detector}. That approach is fine when the discipline is followed, but it offers no runtime safety net
\parencite{netbsd_deadlock_detector}. The same complex interactions that FreeBSD's tooling would flag in development can
therefore persist into production on NetBSD, which can surface caused by reboot-only freezes \parencite{petermann_jails}.

\subsection{Locking Strategies} \label{locking-strategies} % done

FreeBSD follows a fine-grained locking strategy refined across two decades through the SMPng project, which was previously
covered in Section \ref{scalability} \parencite{freebsd_vs_dragonfly}. To revisit, FreeBSD 5.x replaced the \texttt{Giant}
lock with per-subsystem locks; FreeBSD 7.x added ULE and finer lock granularity; FreeBSD 10 and later moved most subsystems to
per-structure locks, reader-writer locks, or lockless algorithms; and FreeBSD 14 brought mature NUMA-aware allocation
\parencite{freebsd_vs_dragonfly}. This transition exposed a significant drawback of fine-grained locking \parencite{freebsd_netperf}.
It improved parallelism and preemption in the network stack but it substantially raised per-packet processing cost through
synchronization \parencite{freebsd_netperf}. The subsequent \textit{netperf} project was formed specifically to recover
parallelism while reducing the overhead that comes with it \parencite{freebsd_netperf}. The added complexity cost is quantifiable
according to \textcite{toy2010}. An evaluation of FreeBSD locking found that replacing \texttt{Giant} with fine-grained locks
in the \texttt{sysctl} subsystem raised the lock-acquisition count by around two orders of magnitude \parencite{toy2010}.
Correct fine-grained locking acquires a lock at every shared-data access point rather than once per syscall
\parencite{toy2010}. In addition, every such acquisition site is a potential ordering bug \parencite{toy2010}.
To keep that overhead in check, FreeBSD also avoids locking where it can. Widely-shared credential structures use a
copy-on-write discipline \parencite{freebsd_smp_lock_strategies}. Many references may share one \texttt{struct ucred},
and modifying it duplicates the structure and updates only the caller's reference \parencite{freebsd_smp_lock_strategies}.
Under fine-grained SMP, a thread consuming a shared credential then needs no lock at all because the copy-on-write guarantee
ensures that no other thread will mutate its copy \parencite{freebsd_smp_lock_strategies}. Production visibility into
whatever contention remains comes from the DTrace \texttt{lockstat} provider \parencite{freebsd_smp_lock_strategies}. It
reports which locks are spinning, how many events per second occur, and which callers are the primary contenders,
all without the overhead of a debug kernel \parencite{freebsd_smp_lock_strategies}.

NetBSD 4.0 used a coarse model of spinlocks combined with a per-CPU interrupt priority level that accumulated over many
years, whose rules had become difficult to use correctly \parencite{feyrer2008}. The \texttt{lockmgr} primitive, the only
available sleep lock, was inefficient and slow \parencite{feyrer2008}. NetBSD 5.0 replaced this wholesale with fine-grained
locking across almost all core subsystems \parencite{feyrer2008}. However, this conversion was not uniform. The
VM system, the allocators, and most filesystem frameworks were converted, but the network stack remained under a single
global lock \parencite{netbsd_smp_networking}. NetBSD minimizes contention through a hierarchy of increasingly lightweight
read-side mechanisms:

\begin{itemize}
	\item \textbf{Passive serialization} (\texttt{pserialize}) enables lock-free reads in soft-interrupt context
	      \parencite{netbsd_locking}.
	\item \textbf{Passive references} (\texttt{psref}) let a CPU acquire and release a reference to a resource cheaply with
	      no inter-processor synchronization during normal use, deferring that synchronization until the resource is pending
	      destruction \parencite{netbsd_locking}.
	\item \textbf{Per-CPU reference counts} (\texttt{localcount}) maintain reference counts per CPU without atomic
	      operations in the common case \parencite{netbsd_locking}.
\end{itemize}

\noindent
A major tradeoff of these mechanisms is cost asymmetry. While reads are nearly free, but destroying a resource requires
expensive inter-processor synchronization to drain all notable references \parencite{netbsd_locking}. To find the hot locks
that are worth refactoring, NetBSD provides \texttt{lockstat()}, which was first introduced in NetBSD 4.0
\parencite{netbsd_lockstat}. It groups adaptive-mutex spin events by lock and call site with counts and durations
\parencite{netbsd_lockstat}. A developer can thereby identify a contender such as \texttt{uvm\_pageqlock} appearing under
\texttt{uvm\_fault\_internal} and decide whether to move to a finer-grained or lockless approach \parencite{netbsd_lockstat}.

It can be observed that systems pursue fine-grained locking. Both also supply production-grade contention profilers with
FreeBSD's DTrace \texttt{lockstat} provider and NetBSD's \texttt{lockstat()}. Their strategies mainly differ in completeness
and emphasis. FreeBSD's conversion is essentially total and reaches even the network stack. It then layers copy-on-write and
lockless techniques on top to claw back the per-acquisition overhead that fine granularity imposes \parencite{freebsd_smp_lock_strategies}.
NetBSD's conversion is more deliberate but partial, which leaves the network stack globally locked
\parencite{netbsd_smp_networking}. Though it reaches for a richer set of read-mostly primitives, namely \texttt{pserialize},
\texttt{psref}, and \texttt{localcount}, which push read-side cost toward zero \parencite{netbsd_locking}.

\subsection{Performance and Overhead} \label{concurrency-performance} % done

FreeBSD's blocking mutex is cheaper than a spin mutex in the uncontested case, which requires only a single atomic
compare-and-swap to acquire or release \parencite{mckusick2014_process}. A contested release is far more expensive because the
kernel must locate the turnstile, lock the turnstile chain, and wake the waiters \parencite{mckusick2014_process}.
Adaptive spinning mitigates this by letting a waiter spin when the holder is running on another CPU instead of
entering the costly block-and-wakeup cycle \parencite{mckusick2014_process}. Locking too finely appeared during the netperf work
\parencite{freebsd_netperf}. The first attempt at fine-grained network-stack locking introduced a lock at every critical
structure, wherein the resulting granularity was so excessive that locking overhead overwhelmed the available
parallelism \parencite{freebsd_netperf}. This was fixed by selectively coarsen and amortizing cost by grouping the
processing of similar items such as packets under a single lock instead of locking per item \parencite{freebsd_netperf}.
FreeBSD also keeps its strictest correctness checks out of the hot path by making them opt-in \parencite{mckusick2014_process}.
Production kernels rely on the locking hierarchy, adaptive mutexes, and turnstile priority propagation,
while \texttt{WITNESS} and \texttt{INVARIANTS} are added only in debug kernels with a worst-case cost exceeding
400\% \parencite{mckusick2014_process}. Two FreeBSD tradeoffs are tunable rather than fixed. The interrupt-driven
reception model imposes a context switch per incoming packet, which can devote the entire CPU to interrupt handling
under heavy loads \parencite{freebsd_polling}. FreeBSD's optional device polling, configurable through \texttt{polling()},
disables interrupt-driven reception in favor of polling at clock ticks and in the idle loop \parencite{freebsd_polling}. This
removes the per-packet switch and lets the operator bound how much CPU goes to device handling \parencite{freebsd_polling}.
The scheduler's time quantum is a second such dial \parencite{mckusick2014_process}. Shorter quanta improve
interactive response but might cause more frequent context switches that flush caches and reduce throughput
\parencite{mckusick2014_process}. ULE therefore sizes quanta dynamically, smaller for interactive threads and
larger for CPU-bound batch threads \parencite{mckusick2014_process}. This is to optimize both rather than fixing
a single compromise.

NetBSD frames its primitives around cost-proportional use that matches the expense of a mechanism to the frequency
and criticality of what it protects:

\begin{itemize}
	\item \texttt{atomic\_ops(3)} and \texttt{membar\_ops(3)} are valid in all three execution contexts and carry only the
	      lowest hardware-level cost \parencite{netbsd_locking}.
	\item Adaptive mutexes carry context-dependent overhead that spin when the holder is on another CPU and sleeping
	      otherwise \parencite{netbsd_locking}.
	\item \texttt{pserialize} and \texttt{psref} provide near-zero-cost reads at the price of expensive write-side
	      inter-processor synchronization \parencite{netbsd_locking}.
	\item \texttt{localcount} avoids atomic operations during normal use but pays an inter-processor cost when draining
	      references \parencite{netbsd_locking}.
\end{itemize}

\noindent
The adaptive mutex itself spins only when the holder runs on a different CPU and sleeps otherwise \parencite{ozaki2015}.
This makes it unsuitable for hardware interrupt handlers, which cannot sleep \parencite{ozaki2015}. Performing sleepable
work such as memory allocation while holding one is explicitly discouraged since it degrades performance and
introduces latency spikes \parencite{ozaki2015}. This modern arrangement directly addresses the hidden costs of the pre-5.0
model \parencite{feyrer2008}. There the \texttt{spl} system masked interrupts on the local CPU for every access to a
shared structure \parencite{feyrer2008}. Its rules are also about which level protected which data were hard to reason
about \parencite{feyrer2008}. Replacing \texttt{spl} with per-object locks made the cost of each protection operation
proportional to the scope it guards, which improved both correctness reasoning and average-case performance \parencite{feyrer2008}.

Both systems adopt adaptive mutexes that spin or sleep according to the state of the holder. Both also treat the strictest correctness
checking as a development-time cost to be excluded from production. Lastly, both learned that lock granularity must be
tuned rather than maximized. FreeBSD's distinctive levers are workload-facing, namely the polling escape from interrupt
livelock and the load-adaptive quantum, both of which suit its server-tuning orientation \parencite{freebsd_polling}.
NetBSD's implementation is the consciously cost-graded ladder of primitives listed above, which lets a developer pay
only for the synchronization strength a given access truly needs \parencite{netbsd_locking}. That ladder is
also why NetBSD's tuning decisions are made largely at the point code is written, whereas FreeBSD's \texttt{polling()}
and dynamic quantum remain knobs an administrator adjusts for a particular load \parencite{freebsd_polling}.

\subsection{Scalability and Multicore Considerations} \label{concurrency-scalability} % done

FreeBSD's multicore scalability is the cumulative product of the twenty-year progression as discussed in Section
\ref{scalability}, which reached scalability beyond 100 cores for many workloads \parencite{freebsd_vs_dragonfly}.
On NUMA hardware, FreeBSD places memory round-robin across nodes to average out access latencies for workloads that
cannot be confined to a single domain \parencite{lameter2013}. This lets it perform reasonably on systems that cannot
interleave cache lines in hardware or firmware \parencite{lameter2013}. Moreover, it prevents any one node from becoming
a bottleneck \parencite{lameter2013}. However, the policy does not optimize for maximum locality, hence workloads
with strong locality affinity may do better under a nearest-node placement \parencite{lameter2013}. The scheduler's topology
awareness is itself a scaling cost at high core counts \parencite{bouron2018}. As noted in Section \ref{resource-performance},
\textcite{bouron2018} found that \texttt{sched\_pickcpu()} may scan all cores up to three times when placing a
waking thread, which consumes 13\% of all CPU cycles on scanning alone in sysbench. The very topology awareness that aids
placement can thus dominate runtime for short-lived, wake-heavy workloads as cores multiply \parencite{bouron2018}. The
workload-dependence established by \textcite{cui2011} further supports this finding. FreeBSD scales better than Linux for
socket-intensive operations but might not perform well for file-descriptor and process-intensive ones, with the shared
bottleneck being synchronization primitives on shared kernel data \parencite{cui2011}.

NetBSD's most recent advancement came with NetBSD 10.0 in March 2024, which delivered substantial SMP gains over the
9.x series \parencite{netbsd_10_0}. The improvements include a faster per-directory red-black tree replacing the global
path-lookup cache, a rewritten page allocator with CPU topology awareness, replacement of global counters with per-CPU
counters, reduced lock contention across the VM system, and improved scheduling for hybrid CPU architectures such
as ARM big.LITTLE \parencite{netbsd_10_0}. Benchmarks showed large gains on compute-intensive and filesystem-bound
multicore workloads \parencite{netbsd_10_0}. The move from global to per-CPU counters directly attacks the bottleneck
identified by \textcite{cui2011}, namely synchronization on shared data and the lock thrashing in which contention at
high core counts degrades rather than improves speedup \parencite{cui2011}. Per-CPU counters remove a global counter
that would otherwise become a single point of contention as cores increase \parencite{cui2011}. NetBSD's scalability
limitation remains its network stack, which is still protected by a single global lock \parencite{netbsd_smp_networking}.
This is attributed to the stack's monolithic structure, its asynchronous upcalls between layers, and its many scattered
global variables \parencite{netbsd_smpnet}. The consequence is that network-intensive workloads cannot spread across
multiple cores simultaneously, which is a hard ceiling for server applications \parencite{netbsd_smp_networking}.
The migration path is set out by \textcite{ozaki2015}, wherein it converts individual drivers and protocol handlers
from \texttt{spl} to adaptive mutexes and \texttt{pserialize}. The  \texttt{spl}-protected code runs in all contexts
but inhibits SMP, while mutex-protected code is SMP-safe in thread context but context-dependent in interrupt context
\parencite{ozaki2015}. Because each protocol layer must be converted, tested, and fitted into the lock hierarchy individually,
the transition stays stable but spans multiple major releases \parencite{ozaki2015}.

FreeBSD has already carried fine-grained locking and NUMA awareness through the entire kernel, which includes the
network stack. It also has reached high single-system core counts while still finding fresh bottlenecks at each
generation \parencite{freebsd_guzik}. Its tuning is deliberate, aimed at locality and at large package-build and microbenchmark
workloads \parencite{freebsd_guzik}. NetBSD has made strong recent progress in its VM, allocator, and scheduler subsystems
\parencite{netbsd_10_0}. It also attacks shared-counter contention directly \parencite{netbsd_10_0}. Its globally-locked network
stack still leaves network-intensive multicore \parencite{netbsd_smp_networking}. Both are converting the kernel to scale
across many cores and the network stack is something FreeBSD has accomplished unlike NetBSD.

\subsection{Summary} \label{concurrency-summary} % done

FreeBSD and NetBSD arrived at similar concurrency foundations, which diverged in terms of maturity, completeness, and
emphasis rather than fundamental design. Both abandoned M:N threading in favor of a 1:1 model that is simpler and faster.
However, the threading model was implemented differently. FreeBSD binds each user thread to a distinct kernel thread,
while NetBSD represents both processes and POSIX threads uniformly as light-weight processes. This serves as a single
abstraction that lets one scheduler apply the same time-sharing and real-time policies to every entity. Both supply nearly
identical synchronization primitives such as adaptive mutexes, reader-writer locks, condition variables, and atomics.
Both also follow a prevention-first stance toward deadlock through consistent lock ordering. FreeBSD's strengths lies
on its depth of implementation and the amount of tooling available, which includes a rich and workload-tunable set of
lock types, mature runtime verifiers in \texttt{WITNESS} and \texttt{INVARIANTS} that convert race conditions and
ordering errors into diagnosable and development-time panics, targeted structural defenses against named hazards,
and a complete fine-grained locking. This is at the cost of complexity and synchronization overhead. Replacing Giant
with fine-grained locks raised lock-acquisition counts in some subsystems by two orders of magnitude, Moreover,
acquisition site is a potential ordering bug. NetBSD's strengths lies in leanness, portability, and minimalism.
This is reflected by its smaller primitive set governed by explicit preconditions and a context-validity matrix,
a cost-graded ladder of read-mostly mechanisms that push read-side cost toward zero, and a portability-first design
that imposes weaker atomic-ordering guarantees requiring explicit memory barriers. It falls short in the absence of
a production-grade runtime deadlock or race detector as it relies on convention, the context matrix, and compile-time
\texttt{LOCKDEBUG}. It also has an incomplete fine-grained conversion that leaves the network stack under a single
global lock.

The tradeoffs are clear across the four comparison dimensions. In terms of correctness, FreeBSD's \texttt{WITNESS}
and \texttt{INVARIANTS} provide a runtime safety net that turns concurrency hazards into immediate panics during
development, whereas NetBSD relies more heavily on convention and its context matrix, which leaves complex
codepaths exposed. This is a gap concretely realized in NetBSD's jails work, where \texttt{secmodel\_eval} interactions
with UVM and scheduling paths under mixed CPU and I/O load produced a full freeze. In terms of performance, both adopt
adaptive mutexes that spin or sleep by holder state and both keep strict correctness checks out of the production hot
path. FreeBSD's distinctive levers are workload-facing administrator knobs like polling and the load-adaptive time
quantum. While NetBSD's contribution is its cost-proportional ladder of primitives chosen at code-writing time. In
terms of scalability, both pursue fine-grained and topology-aware locking, but FreeBSD further refined on it
with its benchmark that reached beyond 100 cores with NUMA awareness across the whole kernel including the network stack.
NetBSD made strong recent gains in 10.0, yet still leaves its globally-locked network stack as a hard ceiling for
network-intensive multicore workloads. In terms of complexity, FreeBSD is consistently heavier, as reflected
from its layered copy-on-write and lockless techniques that claws back the overhead its total conversion imposes.
This is in contrast with NetBSD wherein it stays simpler and more uniform, which makes it less complete defers
much tuning to the moment code is written rather than to runtime.

For a system with many threads accessing shared data, both systems track lock owners and waiters and apply
priority inheritance through turnstiles to prevent inversion, but FreeBSD offers finer control. Developers
get to select among lock types, which allows them to balance priority propagation against sleep tolerance.
Furthermore \texttt{rmlock} propagates priority even for read-mostly data. Lastly, \texttt{WITNESS} catches ordering
violations before deployment. NetBSD handles the same scenario with a smaller set nudging developers toward minimal
synchronization, using \texttt{pserialize} for lock-free reads that instead block writers until all readers exit, and leaning on
the context matrix to foreclose cross-context deadlocks. For a high-performance application such as
a web server or database, FreeBSD is generally the stronger performer since its complete fine-grained conversion
lets network-intensive workloads spread across many cores. Also, its zero-copy \texttt{sendfile()} and VM-optimized pipes,
which outperform local sockets by around 63\%, suit high-volume content serving. In addition, its SysV shared
memory and shared/exclusive locks favor database workloads. Its scaling past 100 cores with NUMA awareness
also directly serves server hardware. NetBSD competes on lighter-weight efficiency and portability and made
multicore progress in 10.0, but its single global network lock prevents network-intensive workloads
from spreading across cores. Overall, FreeBSD trades complexity and overhead for completeness, mature tooling, and
server-scale performance, while NetBSD trades feature completeness for minimalism, portability, and
competitive lightweight efficiency on the subsystems it has already converted.
